% ============================================================
% Chapter 4 – Implementation
% chapters/04-implementation.tex
% ============================================================

\chapter{Implementation}
\label{chap:implementation}

\todo[inline]{Y justifier mes choix techniques:
\begin{itemize}
    \item le formulaire j'ai choisi à l'aide de mon tableau de comparaison.csv
    \item cmt j'ai implémenter les shapes
    \item cmt j'aurai géré les imbrications
\end{itemize}
}

This chapter describes the realisation of the conceptual solution presented above,
focusing on the RML-core use case as a concrete evaluation scenario.

\subsection{Choice of Base Form Library}

From the state-of-the-art comparison (see Table~\ref{tab:comparison}), the
\textit{ULB-Darmstadt shacl-form} was selected as the base library. It offers the
broadest support for the SHACL features required by our use case --- in particular
\texttt{sh:or} and \texttt{sh:xone} rendered as dropdowns --- and its open-source,
TypeScript-based architecture makes it amenable to extension.

\subsection{RML Shape Design}

Rather than using the official RML-Core shapes directly --- which are designed for
validation and produce deeply nested, difficult-to-render forms --- we designed a
shape profile (\texttt{rml-shape.ttl}) oriented towards form
generation.
\% TODO Expliquer ce que j'ai transformer en justifiant par les steps décrits au chapter 3.

This profile retains the semantic correctness of RML while reducing
unnecessary nesting (e.g.\ collapsing \texttt{AbstractLogicalSource} into
\texttt{LogicalSource}) and replacing unsupported constructs such as
\texttt{sh:alternativePath} with \texttt{sh:xone}.

\subsection{Multi-Node and Node Reuse}

\todo[inline]{Décrire concrètement: (1) le wrapper shape ex:Mapping ou une autre solution,
(2) l'utilisation/adaptation de setClassInstanceProvider pour la réutilisation
de nœuds créés pendant la session.}